% Options for packages loaded elsewhere
\PassOptionsToPackage{unicode}{hyperref}
\PassOptionsToPackage{hyphens}{url}
\documentclass[
  a4paper,
]{article}
\usepackage{xcolor}
\usepackage[margin=25mm]{geometry}
\usepackage{amsmath,amssymb}
\setcounter{secnumdepth}{-\maxdimen} % remove section numbering
\usepackage{iftex}
\ifPDFTeX
  \usepackage[T1]{fontenc}
  \usepackage[utf8]{inputenc}
  \usepackage{textcomp} % provide euro and other symbols
\else % if luatex or xetex
  \usepackage{unicode-math} % this also loads fontspec
  \defaultfontfeatures{Scale=MatchLowercase}
  \defaultfontfeatures[\rmfamily]{Ligatures=TeX,Scale=1}
\fi
\usepackage[]{lmodern}
\ifPDFTeX\else
  % xetex/luatex font selection
\fi
% Use upquote if available, for straight quotes in verbatim environments
\IfFileExists{upquote.sty}{\usepackage{upquote}}{}
\IfFileExists{microtype.sty}{% use microtype if available
  \usepackage[]{microtype}
  \UseMicrotypeSet[protrusion]{basicmath} % disable protrusion for tt fonts
}{}
\makeatletter
\@ifundefined{KOMAClassName}{% if non-KOMA class
  \IfFileExists{parskip.sty}{%
    \usepackage{parskip}
  }{% else
    \setlength{\parindent}{0pt}
    \setlength{\parskip}{6pt plus 2pt minus 1pt}}
}{% if KOMA class
  \KOMAoptions{parskip=half}}
\makeatother
\usepackage{longtable,booktabs,array}
\newcounter{none} % for unnumbered tables
\usepackage{calc} % for calculating minipage widths
% Correct order of tables after \paragraph or \subparagraph
\usepackage{etoolbox}
\makeatletter
\patchcmd\longtable{\par}{\if@noskipsec\mbox{}\fi\par}{}{}
\makeatother
% Allow footnotes in longtable head/foot
\IfFileExists{footnotehyper.sty}{\usepackage{footnotehyper}}{\usepackage{footnote}}
\makesavenoteenv{longtable}
\usepackage{graphicx}
\makeatletter
\newsavebox\pandoc@box
\newcommand*\pandocbounded[1]{% scales image to fit in text height/width
  \sbox\pandoc@box{#1}%
  \Gscale@div\@tempa{\textheight}{\dimexpr\ht\pandoc@box+\dp\pandoc@box\relax}%
  \Gscale@div\@tempb{\linewidth}{\wd\pandoc@box}%
  \ifdim\@tempb\p@<\@tempa\p@\let\@tempa\@tempb\fi% select the smaller of both
  \ifdim\@tempa\p@<\p@\scalebox{\@tempa}{\usebox\pandoc@box}%
  \else\usebox{\pandoc@box}%
  \fi%
}
% Set default figure placement to htbp
\def\fps@figure{htbp}
\makeatother
\setlength{\emergencystretch}{3em} % prevent overfull lines
\providecommand{\tightlist}{%
  \setlength{\itemsep}{0pt}\setlength{\parskip}{0pt}}
\usepackage{bookmark}
\IfFileExists{xurl.sty}{\usepackage{xurl}}{} % add URL line breaks if available
\urlstyle{same}
\hypersetup{
  hidelinks,
  pdfcreator={LaTeX via pandoc}}

\author{}
\date{}

\begin{document}

\section{Plane-Decomposition Readout in N-Body Fixed-Generator
Systems}\label{plane-decomposition-readout-in-n-body-fixed-generator-systems}

\subsection{Each rotation plane reads as kinematically isomorphic to the
two-body AB system --- plane-wise validity of progression, frequency,
amplitude, energy, and the zero-order harmonic structure, with position
and observed-acceleration readouts as connection
problems}\label{each-rotation-plane-reads-as-kinematically-isomorphic-to-the-two-body-ab-system-plane-wise-validity-of-progression-frequency-amplitude-energy-and-the-zero-order-harmonic-structure-with-position-and-observed-acceleration-readouts-as-connection-problems}

\textbf{Author:} Noriaki Kihara\\
\textbf{ORCID:} 0009-0004-6753-4020\\
\textbf{Date:} July 21, 2026\\
\textbf{Version DOI:} 10.5281/zenodo.21468960\\
\textbf{Concept DOI:} 10.5281/zenodo.21468959\\
\textbf{Position:} Dimensional Generation Structure series, Paper 4, v1

\begin{center}\rule{0.5\linewidth}{0.5pt}\end{center}

\subsection{Abstract}\label{abstract}

The preceding Wave Information Readout series established, in the closed
two-body AB phase system, position and harmonic readout {[}2{]} and the
acceleration-like readout with the inverse-square law {[}3{]}. Paper 3
{[}5{]} proved that the fixed generator of N-body complete pairwise
relational waves decomposes into a direct sum of rotation planes and a
null space, with the number of rotation modes growing at most linearly.
This paper connects the two series and answers the following question.

\begin{quote}
From which part of an N-body fixed-generator system, and to what extent,
do the readouts established for the two-body AB system hold?
\end{quote}

The key is that under a fixed generator each rotation plane is an
invariant subspace and the motion on it is a single-frequency phase
progression. We formalize this as the plane AB kinematic isomorphism and
establish five mathematical consequences. First, the phase progression
of each plane advances by the constant angle
\(\theta_j=2\arctan(\gamma\sigma_j)\) with constant plane amplitude
(Theorem 1). Second, under the in-plane \(SO(2)\) gauge (the orientation
of each plane is fixed by the generator), the absolute phase is
unreadable while the signed phase progression, frequency, and amplitude
are readable; if reflections are allowed (\(O(2)\)), only the absolute
values of progression and frequency are invariant. This means that Axiom
5 of the basic axiom system (unreadability of individual phases) is
reproduced at the level of derived quantities (Theorem 2). Third, the
total energy decomposes exactly into per-plane and kernel conserved
quantities: \(H=\sum_j H_j+H_{\ker}\) (Theorem 3). Fourth, because the
residual subspace (kernel) carries no phase progression,
progression-based readouts --- position as progressing phase and
harmonic acceleration --- do not hold from the kernel; the projector,
the kernel squared quantity, and the state projection \(P_0X\) are
covariantly defined, while \(K\) alone cannot select a standard internal
basis (the dynamical realization of Consequence 16.3 {[}1{]}, Theorem
4). Whether the static \(P_0X\) can be read as a positional direction is
not tested in this paper. Fifth, the squared closure itself also
decomposes plane-wise and kernel-wise with each term individually
conserved: \(\sum_jQ_j+Q_{\ker}=R^2\) (Theorem 5). Under a fixed
generator, therefore, neither closure nor energy-like quantities move
between planes --- a no-go condition for spontaneous splitting of the
state.

In the numerical experiments (\(N=3\)--\(9\), 32 trials × 720 steps
each, fixed generator, no sequential renormalization) we confirmed:
plane-amplitude drift \(\le1.21\times10^{-13}\); per-step progression
variation \(\le1.89\times10^{-13}\); deviation of the signed mean
progression from \(+2\arctan(\gamma\sigma_j)\)
\(\le4.49\times10^{-15}\); signed-progression change under the in-plane
\(SO(2)\) gauge \(\le1.52\times10^{-13}\); deviation of the orientation
convention \(Kp_j=\sigma_jq_j\) \(\le7.64\times10^{-14}\); deviations of
the permuted frequency and energy sets \(\le3.30\times10^{-15}\);
frequency nondegeneracy (minimum gap \(\ge2.0\times10^{-3}\)). The
second-difference relation \(\Delta^2c_j=-4\sin^2(\theta_j/2)c_j\) held
with absolute deviation \(\le1.27\times10^{-16}\) (machine precision),
and the distance exponent in the continuous-spectrum re-expression, for
all configurations with at least two planes (\(N=4\)--\(9\)), was
\(-1.99999\) to \(-1.99987\). The energy-decomposition identity held to
\(\le2.00\times10^{-15}\), the closure-decomposition identity
\(\sum Q_j+Q_{\ker}-R^2\) to \(\le1.89\times10^{-15}\), the per-plane
closure variation to \(\le1.81\times10^{-13}\), the kernel component was
static to \(\le4.57\times10^{-14}\), the kernel-projector gauge
variation was \(\le7.77\times10^{-16}\), and the angle between
independent kernel bases was at least \(13.6^\circ\).

This is a verification paper: most results were predicted as
consequences of the structure of Paper 3. The substantive novelty is
twofold. First, the asymmetric structure of readability --- on rotation
planes, progression, frequency, amplitude, energy, and the zero-order
harmonic structure hold through the kinematic isomorphism, while
progression-based readouts fail on the residual subspace --- is fixed by
both theorems and numerics. Second, the fixed generator carries a double
conservation register per plane and kernel (\(H_j\) and \(Q_j\)), and
neither quantity moves between planes (the spontaneous-splitting no-go,
Theorem 5). The plane counterpart of the AB same-time two-wave
positional phase readout (\(D_{AB}\)) and the construction of the
observed acceleration of Axiom 17 on the planes are not carried out here
and are stated explicitly as connection problems.

\begin{center}\rule{0.5\linewidth}{0.5pt}\end{center}

\subsection{0. Conclusions}\label{conclusions}

The readability of an N-body fixed-generator system has the following
asymmetric structure.

\[
\boxed{
\begin{aligned}
&\text{Rotation planes (the 2D blocks of }\operatorname{rank}K\text{):}\\
&\qquad\text{progression, frequency, amplitude, energy, and the zero-order harmonic structure all hold (AB kinematic isomorphism).}\\
&\text{Residual subspace (}\ker K\text{):}\\
&\qquad\text{progression-based readouts fail (static); projector, squared quantities, and state projection are covariantly defined.}
\end{aligned}
}
\]

Each rotation plane behaves, as a constant-amplitude single-frequency
two-dimensional rotation, kinematically isomorphic to the two-body AB
system. Phase progression, frequency, amplitude, energy, and the
zero-order harmonic second-order structure lift plane-wise to N bodies
through this isomorphism. By contrast, the plane counterpart of the AB
same-time two-wave positional phase readout (\(D_{AB}\)) and the
construction on the planes of the observed acceleration of Axiom 17
(curvature radius, future phase position, tangential direction) are not
carried out in this paper; they are connection problems.

That under the in-plane \(SO(2)\) gauge the absolute phase is unreadable
while the signed progression is readable shows that Axiom 5
(unreadability of individual phases) {[}1{]} is reproduced not only as
an axiom for elementary waves but also at the level of derived
quantities obtained by spectral decomposition.

The total energy decomposes as

\[
H=\sum_{j=1}^{r}H_j+H_{\ker}
\]

into per-plane and kernel conserved quantities. Moreover, the squared
closure itself decomposes as

\[
X^{\mathsf T}X=\sum_{j=1}^{r}Q_j+Q_{\ker}=R^2,
\qquad
Q_j=(p_j^{\mathsf T}X)^2+(q_j^{\mathsf T}X)^2,
\]

with each term individually conserved (Theorem 5). The fixed generator
carries a double conservation register per plane and kernel; neither
closure nor energy-like quantities move between planes. This is a no-go
condition for spontaneous splitting of the state: splitting requires
sequentially reconstructed generators, nonlinear couplings, or
inter-plane exchange.

The residual subspace is static under the fixed generator and carries no
phase progression, so progression-based readouts (position, harmonic
acceleration) do not hold from the kernel. The projector, the kernel
squared quantity, and the state projection \(P_0X\) are covariantly
defined, and the standard internal basis of the kernel cannot be
selected from \(K\) alone (Consequence 16.3 {[}1{]}). Whether the static
\(P_0X\) can be read as a positional direction is untested.

Spatial adoption (which planes are read as XYZ) is governed by the
adoption conditions of Axiom 16 {[}1{]} and is not the subject of
verification here. What this paper fixes is the validity range of
internal-quantity readout, prior to adoption.

\begin{center}\rule{0.5\linewidth}{0.5pt}\end{center}

\subsection{1. Research Question}\label{research-question}

\subsubsection{1.1 Connecting the two
series}\label{connecting-the-two-series}

The Wave Information Readout series established position and harmonic
readout {[}2{]} and the acceleration-like readout with the
inverse-square law {[}3{]} in the closed two-body AB phase system. All
of these are readouts constructed from the progression of a single phase
relation.

The Dimensional Generation Structure series established the structure of
N-body complete pairwise relational waves. By Paper 3 {[}5{]}, the fixed
generator satisfies

\[
\operatorname{rank}K\le2\min\!\left(N,\left\lfloor\frac{M}{2}\right\rfloor\right)
\]

and decomposes into a direct sum of rotation planes and a null space
(Consequences 16.1, 16.2 {[}1{]}).

The remaining question was the connection of the two series. The AB
readouts read one phase relation; each rotation plane of the N-body
system carries one phase progression. Hence the part of the AB readouts
that can be constructed from a phase progression should lift plane-wise
to N bodies. This paper writes that lift as theorems and verifies it
numerically, and at the same time distinguishes the parts whose lift has
not yet been derived --- the plane counterpart of the same-time two-wave
positional phase readout (\(D_{AB}\)) and the observed-acceleration
projection of Axiom 17 --- as connection problems.

\subsubsection{1.2 Rule on the direction of causation --- plane
coordinates are derived
quantities}\label{rule-on-the-direction-of-causation-plane-coordinates-are-derived-quantities}

Plane coordinates are not the relational waves themselves; they are
derived quantities obtained from the spectral decomposition of the
generator. To call a plane readout a readout, one must first show that
it can be reconstructed covariantly under name permutations from
relational data alone. This is done in Section 3, together with the
strict separation of unreadable quantities (absolute in-plane phase,
internal kernel directions) from readable ones.

No background space is introduced. No distance matrices, no embeddings,
no external observers.

\subsubsection{1.3 What this paper
discriminates}\label{what-this-paper-discriminates}

\begin{enumerate}
\def\labelenumi{\arabic{enumi}.}
\tightlist
\item
  Can each rotation plane be read as a single-frequency rotation
  kinematically isomorphic to the two-body AB system?
\item
  Under the in-plane gauge, what is readable and what is not? Is this
  consistent with Axiom 5?
\item
  Does the energy-like readout decompose into per-plane and kernel
  conserved quantities?
\item
  What can be read from the residual subspace? Do progression-based
  readouts hold there?
\item
  Do the zero-order harmonic structure and the continuous-spectrum
  re-expression exponent hold plane-wise?
\item
  Is all of the above covariant under name permutations?
\end{enumerate}

\begin{center}\rule{0.5\linewidth}{0.5pt}\end{center}

\subsection{2. Classification of Claims}\label{classification-of-claims}

This is a verification paper. All theorems are consequences of the
existing structure (the invariant-subspace decomposition of the fixed
generator), and the numerical experiments verify their implementation.
The axiomatic basis is Basic Axiom System v6 {[}1{]}.

{\def\LTcaptype{none} % do not increment counter
\begin{longtable}[]{@{}
  >{\raggedright\arraybackslash}p{(\linewidth - 4\tabcolsep) * \real{0.3333}}
  >{\raggedright\arraybackslash}p{(\linewidth - 4\tabcolsep) * \real{0.3333}}
  >{\raggedright\arraybackslash}p{(\linewidth - 4\tabcolsep) * \real{0.3333}}@{}}
\toprule\noalign{}
\begin{minipage}[b]{\linewidth}\raggedright
Item
\end{minipage} & \begin{minipage}[b]{\linewidth}\raggedright
Classification
\end{minipage} & \begin{minipage}[b]{\linewidth}\raggedright
Position in this paper
\end{minipage} \\
\midrule\noalign{}
\endhead
\bottomrule\noalign{}
\endlastfoot
N-body complete pairwise relational-wave model & Model definition &
Implementation of Axiom 16 {[}1{]}; inherited from {[}4,5{]} \\
Phase-difference sine generator, fixed-generator dynamics & Working
hypothesis & Working Axiom I of {[}1{]}; identical to {[}4,5{]} \\
Plane coordinates = derived quantities from spectral decomposition &
Definition & Section 3; name-permutation covariance is tested \\
Plane AB kinematic isomorphism (constant progression and amplitude) &
Mathematical consequence & Theorem 1; does not assert independent plane
closure \(Q_j=0\) \\
In-plane \(SO(2)\) gauge and readable quantities & Mathematical
consequence & Theorem 2; under reflection-including \(O(2)\) the signed
progression is not invariant \\
Energy-like decomposition \(H=\sum H_j+H_{\ker}\) & Mathematical
consequence & Theorem 3 \\
Readout restriction on the residual subspace & Mathematical consequence
& Theorem 4; dynamical realization of Consequence 16.3 \\
Plane-wise validity of the zero-order harmonic structure & Mathematical
consequence (corollary of Theorem 1) & Section 5; distinguished from the
observed-acceleration projection of Axiom 17 \\
Plane counterpart of the same-time two-wave positional readout
\(D_{AB}\) & Connection problem & The corresponding plane quantity is
undefined; not constructed here \\
Plane application of the observed acceleration of Axiom 17 & Connection
problem & \(\rho_c\), future phase position, and tangential projection
are defined after construction on adopted planes \\
Closure decomposition \(\sum Q_j+Q_{\ker}=R^2\) and conservation of each
\(Q_j\) & Mathematical consequence & Theorem 5 \\
Spontaneous-splitting no-go under fixed generators & Derived consequence
& Corollary of Theorem 5; splitting requires dynamics beyond the fixed
generator \\
Continuous-spectrum re-expression exponent \(-2\) & Consequence of a
definition & Section 5; follows identically from the cell-width
definition \(\Delta\theta_j=2\pi\theta_{\min}/\theta_j\) \\
Plane version of integer-harmonic closure (the mechanism of {[}3{]}) &
Not established (unexamined) & Plane frequency ratios are generally
non-integer \\
Frequency nondegeneracy & General-position assumption (independent) &
Section 3.3; not derivable from maximal rank; holds numerically in all
224 trials \\
Results of Experiments F--I & Numerical facts & Section 7 \\
Spatial adoption limited to three planes & Rule of Axiom 16 {[}1{]} &
Outside the scope of verification here \\
Behavior of nearly degenerate planes & Out of scope & Treated in the
sequel (allocation and observation selection) \\
First-order readout (velocity feedback) on planes & Out of scope & This
paper is restricted to zero order \\
\end{longtable}
}

\begin{center}\rule{0.5\linewidth}{0.5pt}\end{center}

\subsection{3. Model and Plane
Decomposition}\label{model-and-plane-decomposition}

\subsubsection{3.1 Inheritance of the
model}\label{inheritance-of-the-model}

To the complete pairwise relation set \(\mathcal E_N\)
(\(M=\binom{N}{2}\)) of \(N\) anonymous individuals we assign complex
relational waves \(X\in\mathbb C^M\) with the closure
\(\sum_eX_e^2+(iR)^2=0\) {[}1,4{]}. From the initial phases the
phase-difference sine generator
\(\widetilde K_{ef}=A_{ef}\sin(\theta_f-\theta_e)\) is constructed once,
normalized by its spectral norm, and kept fixed (fixed generator). The
update is the real orthogonal Cayley transform

\[
U=(I-\gamma K)^{-1}(I+\gamma K),\qquad \gamma=\tan(\pi/144),
\]

which preserves the squared closure and the absolute-square sum without
sequential renormalization {[}5,8{]}.

\subsubsection{3.2 Spectral decomposition and plane
coordinates}\label{spectral-decomposition-and-plane-coordinates}

A real antisymmetric matrix \(K\) decomposes under an orthogonal
transformation into two-dimensional rotation blocks and a zero block
{[}7{]}. To each nonzero block \(j\) corresponds an orthonormal pair
\((p_j,q_j)\) and a frequency \(\sigma_j>0\); to the null space
corresponds an orthonormal basis \(\{n_a\}\). The plane coordinate is

\[
\zeta_j(\tau)=p_j^{\mathsf T}\,\Re X_\tau+i\,q_j^{\mathsf T}\,\Re X_\tau
\]

(the same construction runs in parallel for \(\Im X_\tau\)).

\subsubsection{3.3 Well-definedness of the plane
readout}\label{well-definedness-of-the-plane-readout}

The plane readout is grounded as follows. First, \((p_j,q_j,\sigma_j)\)
is determined solely by the spectral decomposition of \(K\) constructed
from relational data (the initial phases), without external axis names.
Second, under a name permutation \(\pi\),
\(K'=P_\pi KP_\pi^{\mathsf T}\), so the frequency set \(\{\sigma_j\}\)
is invariant and the plane bases move covariantly,
\(p_j\mapsto P_\pi p_j\). The readout quantities --- frequencies,
amplitudes, energies --- are therefore permutation invariant (verified
numerically in Experiment F). Third, each plane basis carries a residual
gauge. Since \(K\) itself fixes the orientation of each plane (the
standard form with \(\sigma_j>0\) is fixed by the orientation
\(Kp_j=\sigma_jq_j\)), the residual gauge preserving the standard form
is \(SO(2)\), and by Theorem 2 this freedom does not affect the readout
quantities. Allowing reflections (\(O(2)\)) flips the signed
progression.

When frequencies degenerate, additional freedom remains in the very
separation into planes. Note that the generic rank equality of
Consequence 16.2 {[}1{]} guarantees \(\sigma_j\ne0\);
\(\sigma_j\ne\sigma_k\) (\(j\ne k\)) is a separate proposition. This
paper places frequency nondegeneracy explicitly as an independent
general-position assumption. In all 224 trials of the numerical
experiments (\(N=3\)--\(9\) × 32) all frequencies were distinct. The
proof that the degeneracy set has measure zero is not given here
(Section 8.4). The behavior of near degeneracy is likewise out of scope
(Section 8).

\begin{center}\rule{0.5\linewidth}{0.5pt}\end{center}

\subsection{4. Plane Readout Theorems}\label{plane-readout-theorems}

\begin{quote}
\textbf{Theorem 1 (plane AB kinematic isomorphism)}

Under the fixed-generator dynamics, the coordinate of each rotation
plane satisfies

\[
\zeta_j(\tau+1)=e^{+i\theta_j}\zeta_j(\tau),
\qquad
\theta_j=2\arctan(\gamma\sigma_j).
\]

That is, the phase progression is the constant angle \(\theta_j\) and
the amplitude \(\lvert\zeta_j\rvert\) is constant: each plane is, as a
constant-amplitude single-frequency two-dimensional rotation,
kinematically isomorphic to the rotational kinematics of the two-body AB
system {[}2,3{]}. The progression readout is defined for planes with
nonzero amplitude \(\lvert\zeta_j\rvert>0\) (at \(\zeta_j=0\) the phase
is undefined).
\end{quote}

\textbf{Proof.} On plane \(j\), \(K\) acts as the two-dimensional
antisymmetric block

\[
\begin{pmatrix}0&-\sigma_j\\\sigma_j&0\end{pmatrix},
\]

whose Cayley transform is a two-dimensional rotation by
\(2\arctan(\gamma\sigma_j)\). The blocks are mutually orthogonal and do
not mix, so the plane coordinate rotates at constant angular speed with
conserved amplitude. \(\blacksquare\)

Note that the complex bilinear quadratic quantity
\(Q_j=(p_j^{\mathsf T}X)^2+(q_j^{\mathsf T}X)^2\) within the plane is
also conserved under the real orthogonal update, but \(Q_j=0\) does not
hold. The theorem asserts a kinematic isomorphism, not that each plane
is an independently closed subsystem.

\begin{quote}
\textbf{Theorem 2 (in-plane \(SO(2)\) gauge and readable quantities)}

For nonzero amplitude \(\lvert\zeta_j\rvert>0\), under the in-plane
basis rotation
\((p_j,q_j)\mapsto(p_j\cos\beta+q_j\sin\beta,\,-p_j\sin\beta+q_j\cos\beta)\)
the absolute phase \(\arg\zeta_j\) shifts by \(-\beta\), while the
signed phase progression \(\Delta\varphi_j\), the frequency
\(\theta_j\), and the amplitude \(\lvert\zeta_j\rvert\) are invariant.
Under the \(SO(2)\) gauge, therefore, the uniquely readable quantities
are the signed progression, frequency, amplitude, and phase differences;
the absolute phase is not readable.
\end{quote}

\textbf{Proof.} The basis rotation induces
\(\zeta_j\mapsto e^{-i\beta}\zeta_j\). The ratio
\(\zeta_j(\tau+1)/\zeta_j(\tau)\) and the modulus do not depend on
\(e^{-i\beta}\). \(\blacksquare\)

The residual gauge is \(SO(2)\) because \(K\) itself fixes the
orientation of the plane (Section 3.3). Under the reflection
\((p_j,q_j)\mapsto(p_j,-q_j)\), allowed in \(O(2)\), we get
\(\zeta_j\mapsto\overline{\zeta_j}\) and the signed progression flips,
\(\Delta\varphi_j\mapsto-\Delta\varphi_j\); the invariants are then
\(\lvert\Delta\varphi_j\rvert\) and the absolute value of the frequency.

This shows that the unreadability of individual phases, placed as Axiom
5 {[}1{]} for elementary waves, is automatically reproduced at the level
of the derived plane coordinates. The distinction between readable and
unreadable quantities is preserved as one climbs the readout hierarchy.

\begin{quote}
\textbf{Theorem 3 (energy-like decomposition)}

\[
H(X)=\sum_{j=1}^{r}H_j+H_{\ker},
\qquad
H_j=\lvert p_j^\dagger X\rvert^2+\lvert q_j^\dagger X\rvert^2,
\qquad
H_{\ker}=\lVert P_0X\rVert^2
\]

holds as an identity, and under the fixed-generator dynamics each
\(H_j\) and \(H_{\ker}\) is individually conserved.
\end{quote}

\textbf{Proof.} \(\{p_j,q_j\}\cup\{n_a\}\) is an orthonormal basis of
\(\mathbb R^M\), so the decomposition is Parseval's identity. \(U\)
decomposes into rotations within each plane and the identity on the
kernel, so each component norm is invariant. \(\blacksquare\)

\begin{quote}
\textbf{Theorem 4 (readout restriction on the residual subspace)}

The kernel component \(P_0X_\tau\) does not depend on \(\tau\) (static).
Hence readouts based on nonzero phase progression --- position as
progressing phase and harmonic acceleration --- do not hold from the
residual subspace. The projector \(P_0\), the kernel squared quantity
\(H_{\ker}\), and the state projection \(P_0X\) are all covariantly
defined under name permutations. The standard internal basis of the
kernel cannot be selected from \(K\) alone.
\end{quote}

\textbf{Proof.} If \(Kn=0\) then \(Un=n\), so the kernel component is
invariant. The progression is zero, so phase increments and second
differences vanish identically and progression-based readouts are
undefined. Uniqueness of the projector and non-uniqueness of the
internal basis follow from Consequence 16.3 {[}1{]}. \(P_0X\) is
uniquely determined by the state and moves covariantly,
\(P_0X\mapsto P_\pi P_0X\), under permutations. \(\blacksquare\)

Theorem 4 shows that the uniqueness distinctions stated algebraically by
Consequence 16.3 are realized directly as readout availability in the
dynamics. Whether the static state projection \(P_0X\) can be read as a
``positional direction'' is not tested here. The interpretation of Paper
2, which read a one-dimensional kernel as the static Z-normal direction
{[}4{]}, is a precedent located in this untested region and does not
contradict the theorem.

\begin{quote}
\textbf{Theorem 5 (closure decomposition and splitting no-go)}

The squared closure decomposes as

\[
X^{\mathsf T}X
=\sum_{j=1}^{r}Q_j+Q_{\ker}
=R^2,
\qquad
Q_j=(p_j^{\mathsf T}X)^2+(q_j^{\mathsf T}X)^2,
\qquad
Q_{\ker}=\sum_a(n_a^{\mathsf T}X)^2,
\]

and under the fixed-generator dynamics each \(Q_j\) and \(Q_{\ker}\) is
individually conserved.
\end{quote}

\textbf{Proof.} Since \(\{p_j,q_j\}\cup\{n_a\}\) is an orthonormal basis
of \(\mathbb R^M\), the decomposition is Parseval's identity for the
complex bilinear form. For the in-plane rotation
\(c'=c\cos\theta-d\sin\theta\), \(d'=c\sin\theta+d\cos\theta\), the
identity \((c')^2+(d')^2=c^2+d^2\) holds algebraically also for complex
values \(c,d\); on the kernel the update is the identity.
\(\blacksquare\)

\begin{quote}
\textbf{Corollary (no-go for spontaneous splitting)}

Under a fixed generator, neither the closure quantities \(Q_j\) nor the
energy-like quantities \(H_j\) move between planes or between planes and
the kernel. Spontaneous splitting of the state accompanied by
redistribution among registers therefore cannot occur under
fixed-generator dynamics. Splitting requires sequentially reconstructed
generators, nonlinear couplings, or inter-plane exchange --- dynamics
beyond the fixed generator.
\end{quote}

Theorem 5 shows that the fixed-generator system carries a double
conservation register (\(H_j\) and \(Q_j\)) per plane and kernel, and
serves as the starting point for the sequels on allocation, observation
selection, and state splitting.

\begin{center}\rule{0.5\linewidth}{0.5pt}\end{center}

\subsection{5. The Zero-Order Harmonic Structure on
Planes}\label{the-zero-order-harmonic-structure-on-planes}

By Theorem 1 the plane coordinate rotates by the constant angle
\(\theta_j\), so the real coordinate
\(c_j(\tau)=p_j^{\mathsf T}\Re X_\tau\) satisfies the discrete
second-difference relation

\[
\Delta^2c_j
=-4\sin^2\!\left(\frac{\theta_j}{2}\right)c_j
\]

(the zero-order harmonic second-order structure). This is a direct
corollary of Theorem 1, of the same form as the second-order structure
confirmed for the AB system in the inverse-square paper {[}3{]}.

It must be noted that this is not the observed acceleration of Axiom 17
{[}1{]} itself. The readout
\(\boldsymbol\alpha_{\mathrm{read}}=\rho_c\lvert\omega\rvert^2\hat{\boldsymbol t}_\tau\)
of Axiom 17 requires the curvature radius \(\rho_c\), the future phase
position, the circumferential direction \(\hat{\boldsymbol t}_\tau\),
and the projection of the center-directed compensation onto the
tangential direction; none of these are constructed on the planes in
this paper. Only after Axiom 17 is applied on planes adopted as spatial
directions by Axiom 16 can one speak of an observed acceleration. That
construction is a connection problem. What this section and Experiment G
confirm on all planes is the center-directed zero-order harmonic
second-order structure.

To organize the coefficients across planes we define, relative to the
slowest plane,

\[
n_j=\frac{\theta_j}{\theta_{\min}},
\qquad
\Delta\theta_j=\frac{2\pi}{n_j}=2\pi\frac{\theta_{\min}}{\theta_j}.
\]

An important caution: \(n_j\) is \textbf{generally not an integer}. This
is therefore a different structure from the integer-harmonic closure
\(\omega_n=n\omega_1\) (\(n\in\mathbb Z\setminus\{0\}\)) used in the
inverse-square paper {[}3{]}; it is a re-definition of the cell width on
a continuous spectrum (continuous-spectrum re-expression). Under this
definition, in the small-angle regime,

\[
4\sin^2\!\left(\frac{\theta_j}{2}\right)\simeq\theta_j^2\propto\frac{1}{\Delta\theta_j^2}
\]

follows identically. The exponent \(-2\) established here is thus a
consequence of the definition; the plane version of the integer-harmonic
inverse-square law of {[}3{]} is not claimed (unexamined). Moreover,
this quantity is the normalized harmonic coefficient
\(g_j=4\sin^2(\theta_j/2)\); the plane acceleration amplitude is
\(a_j=\rho_jg_j\), depending on the plane amplitude \(\rho_j\). The
exponent \(-2\) is the exponent of the coefficient, not of the
cross-plane acceleration amplitude.

The second-order structure of this paper is restricted to zero order
(fixed generator, frozen \(\omega\)). The plane application of the
first-order readout (velocity feedback, {[}3{]} v2) is out of scope and
treated in a sequel.

\begin{center}\rule{0.5\linewidth}{0.5pt}\end{center}

\subsection{6. The Boundary of Spatial
Adoption}\label{the-boundary-of-spatial-adoption}

What this paper confirms is the readout validity of the \textbf{internal
quantities} of all \(r\) planes --- progression, frequency, amplitude,
energy, and the zero-order harmonic structure. Which planes are adopted
as spatial directions (XYZ) is governed by the adoption conditions of
Axiom 16 (distinguishability by rotation frequencies plus normal
uniqueness, restricted to \(d=3\)) {[}1{]}, and as Paper 3 showed,
adoption saturates at three directions. The experiments here do not
verify the adoption limit. That internal quantities are readable on all
planes and that spatial adoption is limited to three planes coexist
without contradiction --- this is the two-layer structure of readout.

\begin{center}\rule{0.5\linewidth}{0.5pt}\end{center}

\subsection{7. Numerical Experiments}\label{numerical-experiments}

\subsubsection{7.1 Common conditions}\label{common-conditions}

{\def\LTcaptype{none} % do not increment counter
\begin{longtable}[]{@{}
  >{\raggedright\arraybackslash}p{(\linewidth - 2\tabcolsep) * \real{0.4286}}
  >{\raggedleft\arraybackslash}p{(\linewidth - 2\tabcolsep) * \real{0.5714}}@{}}
\toprule\noalign{}
\begin{minipage}[b]{\linewidth}\raggedright
Item
\end{minipage} & \begin{minipage}[b]{\linewidth}\raggedleft
Value
\end{minipage} \\
\midrule\noalign{}
\endhead
\bottomrule\noalign{}
\endlastfoot
Configurations & \(N=3\)--\(9\) \\
Trials & 32 each \\
Steps & 720 each \\
Dynamics & Fixed generator, Cayley orthogonal update
(\(\gamma=\tan(\pi/144)\)) \\
\(R^2\) / imaginary seed amplitude & 1.0 / 0.35 \\
Random seed & 20260731 \\
Tolerance & \(10^{-10}\) \\
Sequential renormalization / observation damping / absolute background
axes & None \\
\end{longtable}
}

Plane counts and kernel dimensions agreed across all 32 trials per
configuration, reproducing the rank law of Paper 3,
\(r=\min(N,\lfloor M/2\rfloor)\), \(\dim\ker K=M-2r\). No planes were
excluded for low amplitude in any trial.

{\def\LTcaptype{none} % do not increment counter
\begin{longtable}[]{@{}rrrr@{}}
\toprule\noalign{}
\(N\) & \(M\) & Rotation planes \(r\) & \(\dim\ker K\) \\
\midrule\noalign{}
\endhead
\bottomrule\noalign{}
\endlastfoot
3 & 3 & 1 & 1 \\
4 & 6 & 3 & 0 \\
5 & 10 & 5 & 0 \\
6 & 15 & 6 & 3 \\
7 & 21 & 7 & 7 \\
8 & 28 & 8 & 12 \\
9 & 36 & 9 & 18 \\
\end{longtable}
}

\subsubsection{7.2 Experiment F: plane AB kinematic
isomorphism}\label{experiment-f-plane-ab-kinematic-isomorphism}

{\def\LTcaptype{none} % do not increment counter
\begin{longtable}[]{@{}
  >{\raggedright\arraybackslash}p{(\linewidth - 2\tabcolsep) * \real{0.4286}}
  >{\raggedleft\arraybackslash}p{(\linewidth - 2\tabcolsep) * \real{0.5714}}@{}}
\toprule\noalign{}
\begin{minipage}[b]{\linewidth}\raggedright
Quantity
\end{minipage} & \begin{minipage}[b]{\linewidth}\raggedleft
Maximum over all configurations
\end{minipage} \\
\midrule\noalign{}
\endhead
\bottomrule\noalign{}
\endlastfoot
Plane-amplitude drift & \(1.21\times10^{-13}\) \\
Per-step progression variation (single-frequency property) &
\(1.89\times10^{-13}\) \\
Deviation of \textbf{signed} mean progression from
\(+2\arctan(\gamma\sigma_j)\) & \(4.49\times10^{-15}\) \\
Signed-progression change under the in-plane \(SO(2)\) gauge &
\(1.52\times10^{-13}\) \\
Deviation of the orientation convention \((Kp_j)\cdot q_j/\sigma_j-1\) &
\(7.64\times10^{-14}\) \\
Deviation of the permuted frequency set & \(2.00\times10^{-15}\) \\
Deviation of the permuted plane-energy set & \(3.30\times10^{-15}\) \\
\end{longtable}
}

All claims of Theorems 1 and 2 held three orders below tolerance. Under
the in-plane \(SO(2)\) gauge the absolute phase moved while the signed
progression was invariant. In all 224 trials the plane frequencies were
all distinct (minimum frequency gap \(2.0\times10^{-3}\) to
\(5.4\times10^{-2}\)), so the general-position assumption of frequency
nondegeneracy held numerically. The rank law (plane count
\(=\min(N,\lfloor M/2\rfloor)\), kernel dimension \(=M-2r\)) also held
in every trial.

\begin{figure}
\centering
\pandocbounded{\includegraphics[keepaspectratio,alt={Plane AB kinematic isomorphism and kernel staticity}]{figs/plane_ab_isomorphism_v1.png}}
\caption{Plane AB kinematic isomorphism and kernel staticity}
\end{figure}

Figure 1 shows a representative \(N=6\) trial. The phase progressions of
the six planes stay at constant values (order \(10^{-2}\)) while the
change of the kernel component stays at machine precision (order
\(10^{-16}\)): the asymmetry of readability appears as a fourteen-order
separation.

\subsubsection{7.3 Experiment G: plane-wise confirmation of the
zero-order harmonic
structure}\label{experiment-g-plane-wise-confirmation-of-the-zero-order-harmonic-structure}

The absolute deviation of the regression slope of
\(\Delta^2c_j=-4\sin^2(\theta_j/2)c_j\) from the theoretical coefficient
was at most \(1.27\times10^{-16}\) (machine precision) over all
configurations. The relative deviation is amplified up to
\(2.5\times10^{-9}\) on the slowest planes, a conditioning effect of the
measure caused by the smallness of the theoretical coefficient itself;
the acceptance criterion is the absolute deviation.

The exponent in the continuous-spectrum re-expression, for all
configurations with at least two planes (\(N=4\)--\(9\)) and all trials,
was

\[
-1.99999\le\text{exponent}\le-1.99987,
\]

matching \(-2\) (deviation \(\le1.2\times10^{-4}\), a known discrete
effect of \(4\sin^2(\theta/2)\) versus \(\theta^2\)).

\begin{figure}
\centering
\pandocbounded{\includegraphics[keepaspectratio,alt={Inverse-square bookkeeping across planes}]{figs/inverse_square_across_planes_v1.png}}
\caption{Inverse-square bookkeeping across planes}
\end{figure}

Figure 2 shows the nine planes of a representative \(N=9\) trial lying
on the slope \(-2\) line in the continuous-spectrum cell-width display.
This visualizes the definitional consequence of Section 5, not a
verification of integer-harmonic closure.

\subsubsection{7.4 Experiment H: energy-like decomposition and closure
decomposition}\label{experiment-h-energy-like-decomposition-and-closure-decomposition}

{\def\LTcaptype{none} % do not increment counter
\begin{longtable}[]{@{}
  >{\raggedright\arraybackslash}p{(\linewidth - 2\tabcolsep) * \real{0.4286}}
  >{\raggedleft\arraybackslash}p{(\linewidth - 2\tabcolsep) * \real{0.5714}}@{}}
\toprule\noalign{}
\begin{minipage}[b]{\linewidth}\raggedright
Quantity
\end{minipage} & \begin{minipage}[b]{\linewidth}\raggedleft
Maximum over all configurations
\end{minipage} \\
\midrule\noalign{}
\endhead
\bottomrule\noalign{}
\endlastfoot
Per-plane energy \(H_j\) drift & \(2.41\times10^{-13}\) \\
Kernel energy \(H_{\ker}\) drift & \(4.55\times10^{-14}\) \\
Identity \(\sum_jH_j+H_{\ker}-H\) & \(2.00\times10^{-15}\) \\
Per-plane closure \(Q_j\) drift & \(1.81\times10^{-13}\) \\
Kernel closure \(Q_{\ker}\) drift & \(4.03\times10^{-14}\) \\
Closure-decomposition identity \(\sum_jQ_j+Q_{\ker}-R^2\) &
\(1.89\times10^{-15}\) \\
\end{longtable}
}

\begin{figure}
\centering
\pandocbounded{\includegraphics[keepaspectratio,alt={Conservation of the energy-like decomposition}]{figs/plane_energy_decomposition_v1.png}}
\caption{Conservation of the energy-like decomposition}
\end{figure}

Figure 3 summarizes the drifts over all configurations. The total energy
decomposes exactly into per-plane plus kernel individually conserved
quantities. The squared closure \(R^2\) likewise decomposes in the same
basis, with each \(Q_j\) individually conserved (numerical confirmation
of Theorem 5). Under the fixed generator, neither conservation register
moves between planes.

\subsubsection{7.5 Experiment I: readout restriction on the residual
subspace}\label{experiment-i-readout-restriction-on-the-residual-subspace}

{\def\LTcaptype{none} % do not increment counter
\begin{longtable}[]{@{}
  >{\raggedright\arraybackslash}p{(\linewidth - 2\tabcolsep) * \real{0.4286}}
  >{\raggedleft\arraybackslash}p{(\linewidth - 2\tabcolsep) * \real{0.5714}}@{}}
\toprule\noalign{}
\begin{minipage}[b]{\linewidth}\raggedright
Quantity
\end{minipage} & \begin{minipage}[b]{\linewidth}\raggedleft
Value
\end{minipage} \\
\midrule\noalign{}
\endhead
\bottomrule\noalign{}
\endlastfoot
Kernel-component time variation (staticity) & max
\(4.57\times10^{-14}\) \\
Kernel-projector gauge variation (uniqueness) & max
\(7.77\times10^{-16}\) \\
Angle between first directions of independent kernel bases
(non-selectability) & min \(13.6^\circ\) (\(N=6\)) to \(50.9^\circ\)
(\(N=9\)) \\
\end{longtable}
}

The kernel component is static, so progression-based position and
harmonic-acceleration readouts do not hold from the kernel. The
projector and the kernel squared quantity are uniquely readable, and the
state projection \(P_0X\) is covariant.

\begin{figure}
\centering
\pandocbounded{\includegraphics[keepaspectratio,alt={Readout asymmetry of the kernel}]{figs/kernel_readout_asymmetry_v1.png}}
\caption{Readout asymmetry of the kernel}
\end{figure}

Figure 4 contrasts the non-uniqueness of kernel directions (angles
between gauge-equivalent bases) with the uniqueness of the kernel
projector (gauge variation \(\le7.8\times10^{-16}\)).

\begin{center}\rule{0.5\linewidth}{0.5pt}\end{center}

\subsection{8. Scope of Claims}\label{scope-of-claims}

\subsubsection{8.1 Mathematical consequences
proven}\label{mathematical-consequences-proven}

\begin{enumerate}
\def\labelenumi{\arabic{enumi}.}
\tightlist
\item
  Plane AB kinematic isomorphism: each plane is a constant-angle
  \(\theta_j=2\arctan(\gamma\sigma_j)\), constant-amplitude
  single-frequency two-dimensional rotation (Theorem 1).
\item
  Under the in-plane \(SO(2)\) gauge the absolute phase is unreadable
  while signed progression, frequency, and amplitude are readable; under
  \(O(2)\) only absolute values (Theorem 2).
\item
  Energy-like decomposition \(H=\sum H_j+H_{\ker}\) with individual
  conservation (Theorem 3).
\item
  Readout restriction on the residual subspace: progression-based
  readouts fail; projector, squared quantities, and state projection are
  covariant (Theorem 4).
\item
  Closure decomposition \(\sum Q_j+Q_{\ker}=R^2\) with individual
  conservation, and the spontaneous-splitting no-go (Theorem 5).
\end{enumerate}

\subsubsection{8.2 Numerical facts
confirmed}\label{numerical-facts-confirmed}

\begin{enumerate}
\def\labelenumi{\arabic{enumi}.}
\tightlist
\item
  All claims of Theorems 1--5 held for \(N=3\)--\(9\), 32 trials each,
  within tolerance \(10^{-10}\) (measured values mostly below
  \(10^{-13}\)).
\item
  The readout quantities (frequency and energy sets) were permutation
  invariant. The orientation convention \(Kp_j=\sigma_jq_j\) and
  frequency nondegeneracy (an independent general-position assumption)
  held in all trials.
\item
  The zero-order harmonic structure
  \(\Delta^2c_j=-4\sin^2(\theta_j/2)c_j\) held at machine precision, and
  the continuous-spectrum re-expression exponent matched \(-2\)
  (\(N=4\)--\(9\)).
\item
  The rank law and kernel dimensions reproduced the results of Paper 3.
\item
  The individual conservation of \(Q_j\), \(Q_{\ker}\) and the
  closure-decomposition identity held at machine precision (numerical
  confirmation of Theorem 5).
\end{enumerate}

\subsubsection{8.3 Interpretations
adopted}\label{interpretations-adopted}

\begin{enumerate}
\def\labelenumi{\arabic{enumi}.}
\tightlist
\item
  Axiom 5 (unreadability of individual phases) is reproduced not only
  for elementary waves but also at the level of derived plane
  coordinates.
\item
  The readout asymmetry of Consequence 16.3 (projector and squared
  quantities unique; internal directions not) is realized directly as
  readout availability in the dynamics.
\item
  Readout has a two-layer structure: internal-quantity readout holds on
  all planes, while spatial adoption is limited to three planes by Axiom
  16.
\end{enumerate}

\subsubsection{8.4 What this paper does not
fix}\label{what-this-paper-does-not-fix}

\begin{enumerate}
\def\labelenumi{\arabic{enumi}.}
\tightlist
\item
  Behavior of nearly degenerate planes; treated in the sequel
  (allocation and observation selection).
\item
  Proof that the frequency-degeneracy set has measure zero;
  nondegeneracy is placed as an independent general-position assumption.
\item
  Whether the plane version of integer-harmonic closure (the
  inverse-square mechanism of {[}3{]}) holds; the exponent \(-2\) here
  is a definitional consequence of the continuous-spectrum
  re-expression.
\item
  Whether the static state projection \(P_0X\) can be read as a
  positional direction (extension of the Z-normal interpretation of the
  one-dimensional kernel {[}4{]} to multi-dimensional kernels).
\item
  Definition and verification of a plane quantity corresponding to the
  same-time two-wave positional phase readout of the AB system
  (\(D_{AB}=\{\delta,2\pi-\delta\}\)).
\item
  Construction and verification of the observed acceleration of Axiom 17
  (\(\rho_c\), future phase position, tangential projection) on adopted
  planes.
\item
  Plane application of the first-order readout (velocity feedback); this
  paper is restricted to zero order.
\item
  Dynamical verification of the spatial adoption limit (three planes);
  the experiments concern internal-quantity readout only.
\item
  Stability of the plane structure under sequentially reconstructed
  generators.
\item
  Inter-plane interactions (Group B: exchange scattering).
\end{enumerate}

\begin{center}\rule{0.5\linewidth}{0.5pt}\end{center}

\subsection{9. Reproducibility}\label{reproducibility}

The numerical experiments are based on the following programs and
outputs.

\subsubsection{Programs}\label{programs}

\begin{itemize}
\tightlist
\item
  \texttt{dimension\_generation\_structure/contrast\_experiments\_v1/run\_nbody\_plane\_decomposition\_readout\_preliminary\_v1.py}
  (Experiments F--I)
\item
  \texttt{dimension\_generation\_structure/contrast\_experiments\_v1/make\_nbody\_plane\_decomposition\_figures\_v1.py}
  (figures)
\end{itemize}

\subsubsection{Main outputs}\label{main-outputs}

\begin{itemize}
\tightlist
\item
  \texttt{plane\_decomposition\_readout\_result\_v1.json}
\item
  \texttt{plane\_decomposition\_trials\_v1.csv}
\item
  \texttt{plane\_decomposition\_readout\_report\_v1.md}
\item
  Four figures (\texttt{plane\_ab\_isomorphism\_v1.png} etc.)
\end{itemize}

Stored under

\texttt{dimension\_generation\_structure/figs/}.

\begin{center}\rule{0.5\linewidth}{0.5pt}\end{center}

\subsection{10. Final Conclusions}\label{final-conclusions}

The N-body fixed generator decomposes into independent single-frequency
rotation planes and carries per-plane and kernel conservation registers.
Each rotation plane is, as a constant-amplitude single-frequency
two-dimensional rotation, kinematically isomorphic to the rotational
kinematics of AB, lifting phase progression, frequency, amplitude,
energy, and the zero-order harmonic second-order structure. The plane
counterpart of the same-time two-wave positional readout (\(D_{AB}\))
and the construction on planes of the observed-acceleration projection
of Axiom 17 remain as connection problems.

The total energy decomposes exactly into per-plane and kernel conserved
quantities. The residual subspace carries no phase progression, so
progression-based position and harmonic-acceleration readouts do not
hold there; the projector, kernel squared quantity, and state projection
\(P_0X\) are covariantly defined, and whether the static \(P_0X\) reads
as a positional direction is untested.

The conclusions of this paper are therefore:

\[
\boxed{
\begin{aligned}
&\text{the N-body fixed generator decomposes into independent single-frequency rotation planes,}\\
&\text{carries a double conservation register (}H_j\text{ and }Q_j\text{) per plane and kernel,}\\
&\text{and progression-based readouts fail on the residual subspace.}
\end{aligned}
}
\]

This asymmetric readability shows that the uniqueness distinctions given
algebraically by Consequence 16.3 become readout availability in the
dynamics. The no-go of Theorem 5 --- under a fixed generator neither
closure nor energy-like quantities move between planes --- means that
spontaneous splitting of the state requires dynamics beyond the fixed
generator (sequential reconstruction, nonlinear coupling, or inter-plane
exchange). The double registers \(\{H_j,Q_j\}\) and
\(\{H_{\ker},Q_{\ker}\}\) are the foundation for the sequels on N-body
observation selection (weak C readout, strong D observation) and state
splitting.

\begin{center}\rule{0.5\linewidth}{0.5pt}\end{center}

\subsection{References}\label{references}

\subsubsection{Self-citations}\label{self-citations}

\begin{enumerate}
\def\labelenumi{\arabic{enumi}.}
\tightlist
\item
  Noriaki Kihara, ``Anonymous Equal-Amplitude Composite-Wave Model:
  Basic Axiom System v6,'' Zenodo, 2026. Version DOI:
  \href{https://doi.org/10.5281/zenodo.21465984}{10.5281/zenodo.21465984},
  Concept DOI:
  \href{https://doi.org/10.5281/zenodo.21315735}{10.5281/zenodo.21315735}.
\item
  Noriaki Kihara, ``Harmonic Readout and c=1 Area-Sweep Preliminary
  Experiments in the Closed AB Two-Body Phase System, v4,'' Zenodo,
  2026. Version DOI:
  \href{https://doi.org/10.5281/zenodo.21374317}{10.5281/zenodo.21374317},
  Concept DOI:
  \href{https://doi.org/10.5281/zenodo.21318696}{10.5281/zenodo.21318696}.
\item
  Noriaki Kihara, ``An Inverse-Square Law from a Future-Phase-Position
  Acceleration Map and Harmonic Closure in a Closed Two-Body AB Phase
  System, v4,'' Zenodo, 2026. Version DOI:
  \href{https://doi.org/10.5281/zenodo.21468270}{10.5281/zenodo.21468270},
  Concept DOI:
  \href{https://doi.org/10.5281/zenodo.21441081}{10.5281/zenodo.21441081}.
\item
  Noriaki Kihara, ``Three XYZ Directions Read from Complete Pairwise
  Relational Waves: Growth of Internal Relational Directions and
  Saturation of Spatial-Direction Readout in Closed AB, ABC, and ABCD
  Systems,'' Zenodo, 2026. Version DOI:
  \href{https://doi.org/10.5281/zenodo.21454790}{10.5281/zenodo.21454790},
  Concept DOI:
  \href{https://doi.org/10.5281/zenodo.21454789}{10.5281/zenodo.21454789}.
\item
  Noriaki Kihara, ``Linear Upper Bound on Generator Rank and
  Three-Direction Saturation of Spatial-Direction Readout in N-Body
  Complete Pairwise Relational Waves,'' Zenodo, 2026. Version DOI:
  \href{https://doi.org/10.5281/zenodo.21465899}{10.5281/zenodo.21465899},
  Concept DOI:
  \href{https://doi.org/10.5281/zenodo.21465898}{10.5281/zenodo.21465898}.
\end{enumerate}

\subsubsection{External references}\label{external-references}

\begin{enumerate}
\def\labelenumi{\arabic{enumi}.}
\setcounter{enumi}{5}
\tightlist
\item
  Hassler Whitney, ``Congruent Graphs and the Connectivity of Graphs,''
  \emph{American Journal of Mathematics}, 54(1), 150--168, 1932. DOI:
  \href{https://doi.org/10.2307/2371086}{10.2307/2371086}.
\item
  Milan Vujivčić, Fedor Herbut, and Gradimir Vujivčić, ``Canonical Form
  for Matrices Under Unitary Congruence Transformations. I:
  Conjugate-Normal Matrices,'' \emph{SIAM Journal on Applied
  Mathematics}, 23(2), 225--238, 1972. DOI:
  \href{https://doi.org/10.1137/0123025}{10.1137/0123025}.
\item
  Fasma Diele, Luciano Lopez, and R. Peluso, ``The Cayley Transform in
  the Numerical Solution of Unitary Differential Systems,''
  \emph{Advances in Computational Mathematics}, 8(4), 317--334, 1998.
  DOI:
  \href{https://doi.org/10.1023/A:1018908700358}{10.1023/A:1018908700358}.
\end{enumerate}

\end{document}
